% Generated draft: requires mathematical and source review.
\documentclass[11pt]{article}
\usepackage[margin=25mm]{geometry}
\usepackage{fontspec}
\setmainfont{Latin Modern Roman}
\usepackage{amsmath,amssymb,mathtools,unicode-math}
\setmathfont{Latin Modern Math}
\usepackage{enumitem,needspace,xurl,hyperref}
\setlength{\parindent}{0pt}
\setlength{\parskip}{5pt}
\setlength{\emergencystretch}{3em}
\newcommand{\N}{\mathbb{N}}
\newcommand{\Z}{\mathbb{Z}}
\newcommand{\Q}{\mathbb{Q}}
\newcommand{\R}{\mathbb{R}}
\newcommand{\C}{\mathbb{C}}
\newcommand{\F}{\mathbb{F}}
\newcommand{\A}{\mathbb{A}}
\newcommand{\PP}{\mathbb{P}}
\newcommand{\E}{\mathbb{E}}
\newcommand{\eps}{\varepsilon}
\newcommand{\dd}{\,\mathrm{d}}
\begin{document}
\section*{UnsolvedMath ID: 2000014}
\textbf{Source problem code:} catalogue entry
\subsection{Definitions and mathematical statement}
Ball's Problem 14 is to develop criteria for dynamic stability and instability of equilibria in nonlinear elasticity. The source gives no particular elastic evolution equation, constitutive law, domain, load, boundary condition, solution class, or stability topology, so it is a research program rather than one quantified conjecture.

For clarity, one possible precise stability framework is the following. After specifying an elastic evolution with state space $(\mathcal X,d)$, let $z(t;z_0)$ be the solution with initial state $z_0$, and let an equilibrium $z_*$ satisfy $z(t;z_*)=z_*$. In a classical displacement-velocity formulation the state would include both the deformation and its velocity, with equilibrium velocity zero. Lyapunov stability would mean
\[
\forall\varepsilon>0\ \exists\delta>0\ \forall z_0\in\mathcal X:\quad
 d(z_0,z_*)<\delta\Longrightarrow
 d(z(t;z_0),z_*)<\varepsilon
\]
for all nonnegative times on which the specified solution is defined, together with whatever existence requirement the chosen stability notion includes. Instability in this sense is the negation, witnessed by an $\varepsilon_0>0$, arbitrarily close admissible initial states, and corresponding times of departure by at least $\varepsilon_0$. This conditional framework does not select Lyapunov stability over asymptotic, spectral, weak-solution, or orbital stability, and does not treat loss of existence as instability without a stated convention.

The requested criteria should connect the chosen elastic constitutive, equilibrium, and boundary data with the intended dynamic behavior. The supplied summary reports energy-coercivity and linearization-based stability/instability criteria in many settings, citing Slemrod and Pottinger--Slemrod, but says a complete general set of criteria for finite nonlinear elasticity remains ``somewhat open.'' No exact partial theorem or universally applicable criterion is reproduced.
\subsection{Short English statement}
The source asks for conditions that predict whether elastic equilibria remain close to their initial state under small dynamical perturbations or become unstable. The evolution, admissible perturbations, and meaning of stability must be fixed before a single rigorous theorem can be stated.
\subsection{Sources}
\begin{itemize}
\item Local UnsolvedMath registry, permanent ID 2000014, exact source code catalogue entry, ``Some Open Problems in Elasticity --- Dynamic stability''; status partially solved, classification PARTIAL-PROGRESS.
\item Ball, elasticity problems (2002), Problem 14, \url{https://people.maths.ox.ac.uk/ball/Articles%20in%20Conference%20Proceedings%20and%20Books/JMB%202002%20re%20Marsden%2060th.pdf}. The record identifies PDF-text extraction from a public author-hosted PDF, accessed July 29, 2026, and says the source presented this as open.
\item The supplied research summary names Slemrod and Pottinger--Slemrod for energy-coercivity and linearization-based criteria in various settings. No titles, dates, precise hypotheses, or theorem statements are supplied. Its assessment that a complete general theory remains somewhat open is retained as a supplied qualification.
\end{itemize}
\end{document}
